\section{Experiments}
This section outlines the experimental framework and evaluation methodology used to validate the system.
\subsection{Training of the model and exporting to ONNX format}
In this project, the YOLO11s model was used, to deliver good performance with enough speed.
The YOLO model is trained on the dataset using two RTX 3090 GPUs inside the UniPD DEI cluster.
The training dataset has been split into 3 smaller subsets such that:
\begin{itemize}
    \item Each subset contains images (and corresponding labels) which are different than the other datasets (preventing data leakage);
    \item Each subset has a maximum size which allows each one of them to be loaded into RAM.
\end{itemize}
As a result, the training process is divided into three sequential steps: each step loads the model weights from the 
previous step and continues training on the next subset. This allows for a faster and more fault-tolerant training process 
with a negligible penalty on accuracy. The full script is available at \texttt{scripts/yolo\_training.py}. \\ \newline
The model detects a bounding box enclosing the suit and the rank on the corners of a poker card. This means that at most 2 bounding 
boxes can be found for the same card, which makes it easier to find in case of partial occlusions. The bounding boxes are then 
classified as the suit and the rank they enclose. \\ \newline
Since the inference has to be performed in a C++ program, the model has to be exported into a ONNX format and imported 
into the program using a library like ONNX Runtime.
\subsection{Inference}
The inference is subdivided into three sections:
\begin{itemize}
    \item Pre-processing of the image: the image is formatted as a valid input of the Yolo11s model and the ONNX Runtime session. 
    Letterbox padding of the input images is enabled by default, since YOLO models use it for the training phase, but the option 
    can also be disabled when calling the function used for the inference;
    \item Inference of the imported YOLO11s model, using a ONNX Runtime session. If the model used is static, YOLO11 generates 84000 detections;
    \item Post-processing of the results. Each detection result consists of 4 values defining the object's bounding box (position and size), and 
    a confidence score for every possible class which the object might correspond to. Non-maxima suppression is used to keep only 1 bounding box 
    for each detected object.
\end{itemize}
In order to handle detection of small cards in big images, a sliding window approach is used: the image is subdivided into smaller, slightly 
overlapping tiles, with size equal to the image size used for training the model (640px). The overhead introduced by this method is mitigated 
by using multi-threading, where each tile is processed by a separate thread in parallel.
\subsection{Performance metrics}
To quantitatively evaluate the performance of the system, different metrics were used:
\begin{itemize}
    \item Intersection over Union (IoU): ratio between the intersection of the predicted and actual
    bounding boxes and the total area covered by both boxes. It is used to assess how well a predicted bounding box
    aligns with the ground truth. To determine whether a detection is correct, a threshold of 0.5 is typically set (
    if $IoU > 0.5$ the detection is considered a True Positive, while if $IoU < 0.5$ the detection is considered a False Positive);
    \[ IoU = \frac{\text{Area of Intersection}}{\text{Area of Union}} \]
    \item True Positives (TP): occur when the model correctly identifies the presence of a target. In our implementation, a detection is counted as a TP 
    only if the predicted class matches the ground truth and the Intersection over Union (IoU) is greater than or equal to 0.5;
    \item False Positives (FP): occur when the model predicts the presence of a target that is not actually present;
    \item False Negatives (FN): occur when the model fails to detect a target that is actually present;
    \item Precision: measures the proportion of correctly predicted bounding boxes among all detections;
    \[ Precision = \frac{TP}{TP + FP} \]
    \item Recall: measures the proportion of correctly predicted bounding boxes among all ground truth boxes;
    \[ Recall = \frac{TP}{TP + FN} \]
    \item F1-Score: harmonic mean of Precision and Recall used to evaluate how well a classification model performs.
    \[  F1 = 2 \cdot \frac{\text{Precision} \cdot \text{Recall}}{\text{Precision} + \text{Recall}} \]
\end{itemize}